% GENERATED FILE. Edit canonical lessons, then run npm run generate:books.
% canonical-source-hash: fnv1a64:4b09fae65cf2f1cb
% canonical-lessons: BN-C393-dijel, BN-C393-haoya, BN-C393-pangchar, BN-C393-phurano, BN-C393-chalu-kora

\chapter{Fuel and Flat Tyres}
\label{ch:bn-fuel-and-flat-tyres}
\hlchaptermodality{393}

\begin{chapteropening}
\textbf{By the end of this chapter:} \emph{I can say we ran out of fuel, a tyre was flat and the engine would start.}

Complete the last lesson of chapter 393: I can say we ran out of fuel, a tyre was flat and the engine would start.
\end{chapteropening}

\section[\texorpdfstring{\bn{ডিজেল} (ḍijel)}{ডিজেল (ḍijel)}]{\bn{ডিজেল} (ḍijel) — diesel}

\label{lesson:BN-C393-dijel}



\begin{quote}
Before the new one: say the Bengali for a signboard, then the Bengali for to keep close to.
\end{quote}

\subsection*{You'll want to know: \bn{ডিজেল}}
\textbf{\bn{ডিজেল}} — \emph{ḍijel} — \textquotedblleft{}diesel\textquotedblright{}.

\begin{quote}
\bn{ট্রাকগুলো} \bn{ডিজেলে} \bn{চলে।}
\end{quote}

\begin{quote}
\emph{ṭrākgulo ḍijele chôle.}
\end{quote}

\begin{quote}
Lorries run on diesel.
\end{quote}

\begin{cousinweb}
\textbf{\bn{ডিজেল}} (\emph{ḍijel}) is borrowed from English \emph{diesel}, named after the engineer Rudolf Diesel; the English z sound is said with \bn{জ}.
\end{cousinweb}

\begin{grammarlens}[title={What you've built}]
One more word for getting around.
\end{grammarlens}

\subsection*{Your turn}
\begin{itemize}
\raggedright
  \item \emph{Say it:} \emph{ḍijel}
  \item \emph{Say it:} \emph{ḍijel}, once more
  \item \emph{Say it:} \emph{ṭrākgulo ḍijele chôle.} — Lorries run on diesel.
  \item \emph{Recall:} say the Bengali for a signboard, then the Bengali for to keep close to, then say \emph{ḍijel} again
\end{itemize}

\begin{tcolorbox}[breakable,colback=teal!4,colframe=teal!35!black,title={Before you move on}]
What does \bn{ডিজেল} mean? (\textquotedblleft{}Diesel\textquotedblright{}.) Say it once more.
\end{tcolorbox}
\section[\texorpdfstring{\bn{হাওয়া} (hāoyā)}{হাওয়া (hāoyā)}]{\bn{হাওয়া} (hāoyā) — air; a breeze}

\label{lesson:BN-C393-haoya}



\begin{quote}
Before the new one: say the Bengali for to keep close to, then the Bengali for diesel.
\end{quote}

\subsection*{You'll want to know: \bn{হাওয়া}}
\textbf{\bn{হাওয়া}} — \emph{hāoyā} — \textquotedblleft{}air; a breeze\textquotedblright{}.

\begin{quote}
\bn{পেট্রোল} \bn{পাম্পে} \bn{আমরা} \bn{চাকায়} \bn{হাওয়া} \bn{ভরলাম।}
\end{quote}

\begin{quote}
\emph{peṭrol pāmpe āmrā chākāy hāoyā bhôrlām.}
\end{quote}

\begin{quote}
At the petrol station we put air in the tyres.
\end{quote}

Also the breeze you feel: \bn{বাতাস} is the wind; \bn{হাওয়া} is the air moving, or the air in a tyre.

\begin{cousinweb}
\textbf{\bn{হাওয়া}} (\emph{hāoyā}) is usually traced to Arabic \emph{hawāʾ}, \textquotedblleft{}air\textquotedblright{}, which came into Bengali through Persian.
\end{cousinweb}

\begin{grammarlens}[title={What you've built}]
One more word for getting around.
\end{grammarlens}

\subsection*{Your turn}
\begin{itemize}
\raggedright
  \item \emph{Say it:} \emph{hāoyā}
  \item \emph{Say it:} \emph{hāoyā}, once more
  \item \emph{Say it:} \emph{peṭrol pāmpe āmrā chākāy hāoyā bhôrlām.} — At the petrol station we put air in the tyres.
  \item \emph{Recall:} say the Bengali for to keep close to, then the Bengali for diesel, then say \emph{hāoyā} again
\end{itemize}

\begin{tcolorbox}[breakable,colback=teal!4,colframe=teal!35!black,title={Before you move on}]
What does \bn{হাওয়া} mean? (\textquotedblleft{}Air; a breeze\textquotedblright{}.) Say it once more.
\end{tcolorbox}
\section[\texorpdfstring{\bn{পাংচার} (pāngchār)}{পাংচার (pāngchār)}]{\bn{পাংচার} (pāngchār) — a puncture}

\label{lesson:BN-C393-pangchar}



\begin{quote}
Before the new one: say the Bengali for diesel, then the Bengali for air.
\end{quote}

\subsection*{You'll want to know: \bn{পাংচার}}
\textbf{\bn{পাংচার}} — \emph{pāngchār} — \textquotedblleft{}a puncture\textquotedblright{}.

\begin{quote}
\bn{পথে} \bn{সাইকেলের} \bn{চাকা} \bn{পাংচার} \bn{হল।}
\end{quote}

\begin{quote}
\emph{pôthe sāikeler chākā pāngchār hôlô.}
\end{quote}

\begin{quote}
On the way the bicycle got a puncture.
\end{quote}

\bn{পাংচার} \bn{হওয়া}, to get a puncture; the repair stall by the road is a \bn{পাংচারের} \bn{দোকান}.

\begin{cousinweb}
\textbf{\bn{পাংচার}} (\emph{pāngchār}) is borrowed from English \emph{puncture}, with \bn{ং} for the English n before ch.
\end{cousinweb}

\begin{grammarlens}[title={What you've built}]
One more word for getting around.
\end{grammarlens}

\subsection*{Your turn}
\begin{itemize}
\raggedright
  \item \emph{Say it:} \emph{pāngchār}
  \item \emph{Say it:} \emph{pāngchār}, once more
  \item \emph{Say it:} \emph{pôthe sāikeler chākā pāngchār hôlô.} — On the way the bicycle got a puncture.
  \item \emph{Recall:} say the Bengali for diesel, then the Bengali for air, then say \emph{pāngchār} again
\end{itemize}

\begin{tcolorbox}[breakable,colback=teal!4,colframe=teal!35!black,title={Before you move on}]
What does \bn{পাংচার} mean? (\textquotedblleft{}A puncture\textquotedblright{}.) Say it once more.
\end{tcolorbox}
\section[\texorpdfstring{\bn{ফুরানো} (phurāno)}{ফুরানো (phurāno)}]{\bn{ফুরানো} (phurāno) — to run out; to be used up}

\label{lesson:BN-C393-phurano}



\begin{quote}
Before the new one: say the Bengali for air, then the Bengali for a puncture.
\end{quote}

\subsection*{You'll want to know: \bn{ফুরানো}}
\textbf{\bn{ফুরানো}} — \emph{phurāno} — \textquotedblleft{}to run out; to be used up\textquotedblright{}.

\begin{quote}
\bn{পথে} \bn{গাড়ির} \bn{পেট্রোল} \bn{ফুরিয়ে} \bn{গেল।}
\end{quote}

\begin{quote}
\emph{pôthe gāṛir peṭrol phuriye gelô.}
\end{quote}

\begin{quote}
The car ran out of petrol on the way.
\end{quote}

Past (he/she …): \bn{সেই} \bn{পেট্রোল} \bn{ফুরিয়ে} \bn{গেল}, \emph{sei peṭrol phuriye gelô}.

The thing that runs out is the subject: \bn{টাকা} \bn{ফুরিয়ে} \bn{গেল}, the money ran out.

\begin{cousinweb}
\textbf{\bn{ফুরানো}} (\emph{phurāno}) is an everyday Bengali verb whose origin is not settled; CDIAL gives no Bengali form that matches it.
\end{cousinweb}

\begin{grammarlens}[title={What you've built}]
One more word for getting around.
\end{grammarlens}

\subsection*{Your turn}
\begin{itemize}
\raggedright
  \item \emph{Say it:} \emph{phurāno}
  \item \emph{Say it:} \emph{phurāno}, once more
  \item \emph{Say it:} \emph{pôthe gāṛir peṭrol phuriye gelô.} — The car ran out of petrol on the way.
  \item \emph{Recall:} say the Bengali for air, then the Bengali for a puncture, then say \emph{phurāno} again
\end{itemize}

\begin{tcolorbox}[breakable,colback=teal!4,colframe=teal!35!black,title={Before you move on}]
What does \bn{ফুরানো} mean? (\textquotedblleft{}To run out; to be used up\textquotedblright{}.) Say it once more.
\end{tcolorbox}
\section[\texorpdfstring{\bn{চালু} \bn{করা} (chālu kôrā)}{চালু করা (chālu kôrā)}]{\bn{চালু} \bn{করা} (chālu kôrā) — to start (an engine); to switch on}

\label{lesson:BN-C393-chalu-kora}



\begin{quote}
Before the new one: say the Bengali for a puncture, then the Bengali for to run out.
\end{quote}

\subsection*{You'll want to know: \bn{চালু} \bn{করা}}
\textbf{\bn{চালু} \bn{করা}} — \emph{chālu kôrā} — \textquotedblleft{}to start (an engine); to switch on\textquotedblright{}.

\begin{quote}
\bn{চালক} \bn{ইঞ্জিন} \bn{চালু} \bn{করলেন।}
\end{quote}

\begin{quote}
\emph{chālôk injin chālu kôrlen.}
\end{quote}

\begin{quote}
The driver started the engine.
\end{quote}

Past (he/she …): \bn{তিনি} \bn{পাখা} \bn{চালু} \bn{করলেন}, \emph{tini pākhā chālu kôrlen}.

\bn{চালু} on its own is running, in use: \bn{দোকানটা} \bn{চালু}, the shop is open for business.

\begin{cousinweb}
\textbf{\bn{চালু} \bn{করা}} (\emph{chālu kôrā}) joins \bn{চালু}, \textquotedblleft{}going, in working order\textquotedblright{}, probably a Bengali formation on \bn{চাল}, \textquotedblleft{}gait, motion\textquotedblright{}, which CDIAL lists under the reconstructed \emph{calyā} (CDIAL 4722), to \bn{করা}, to do.
\end{cousinweb}

\begin{grammarlens}[title={What you've built}]
One more word for getting around.
\end{grammarlens}

\subsection*{Your turn}
\begin{itemize}
\raggedright
  \item \emph{Say it:} \emph{chālu kôrā}
  \item \emph{Say it:} \emph{chālu kôrā}, once more
  \item \emph{Say it:} \emph{chālôk injin chālu kôrlen.} — The driver started the engine.
  \item \emph{Recall:} say the Bengali for a puncture, then the Bengali for to run out, then say \emph{chālu kôrā} again
\end{itemize}

\begin{tcolorbox}[breakable,colback=teal!4,colframe=teal!35!black,title={Before you move on}]
What does \bn{চালু} \bn{করা} mean? (\textquotedblleft{}To start (an engine); to switch on\textquotedblright{}.) Say it once more.
\end{tcolorbox}
